\subsection{Sprint 2}

No início desta \textit{sprint}, identifiquei a necessidade de refatorar a estrutura do \textit{backend} para seguir as boas práticas de desenvolvimento. Solicitei à \ac{ages} II Luana Thurow que realizasse essa refatoração, incluindo a padronização de nomenclaturas de variáveis, organização de pastas e estrutura de arquivos. O trabalho realizado foi excelente, resultando em um repositório bem organizado que facilitou significativamente o desenvolvimento subsequente.

Minha principal responsabilidade nesta \textit{sprint} era testar o arquivo que o \textit{stakeholder} enviaria, para validar a viabilidade de continuar com um projeto genérico e escalável ou focar em um mapa fixo como prova de conceito. Apesar de ter solicitado o arquivo aos \ac{ages} IV em diversas ocasiões, o \textit{stakeholder} não o enviou a tempo de realizarmos os testes necessários.

Enquanto aguardava o arquivo, acompanhei de perto o desenvolvimento das tarefas dos \ac{ages} II. Conseguimos implementar logo no início da \textit{sprint} a funcionalidade da barra de pesquisa que havia ficado pendente, permitindo que a equipe focasse em tarefas mais complexas. Foi definido que dois \ac{ages} II trabalhariam em conjunto no desenvolvimento de um \textit{endpoint} para popular o banco de dados através de planilhas \ac{csv} --- funcionalidade fundamental para garantir a escalabilidade do projeto para diferentes feiras.

Durante as aulas, em conjunto com os outros \ac{ages} III, confirmamos que, considerando que o projeto rodaria em \textit{hardware} sem acesso à internet durante sua execução, o uso da \ac{aws} não seria necessário e poderia até complicar o fluxo de desenvolvimento, especialmente considerando o tamanho reduzido da equipe de \ac{ages} I e II.

No \textit{code freeze}, ao revisar o código da funcionalidade de importação \ac{csv}, identifiquei um problema crítico: a implementação esperava uma planilha contendo IDs e relacionamentos nas colunas, o que seria inviável para o uso pretendido. Após conversar com os \ac{ages} II responsáveis, descobri que a descrição da tarefa no quadro estava mal formulada e continha informações incorretas. Comuniquei à equipe que esta funcionalidade entraria como débito técnico e que precisaríamos repensar sua implementação.

A apresentação para o \textit{stakeholder} foi desafiadora, pois, sem avanços significativos no \textit{frontend}, os \ac{ages} IV demonstraram apenas as implementações dos \textit{endpoints} através do Postman, o que deixou o Igor visivelmente desapontado. Durante a retrospectiva, discutimos a importância de descrições claras e precisas nas tarefas do quadro, e os \ac{ages} IV se comprometeram a melhorar esse aspecto para as próximas \textit{sprints}.

Embora não tenha contribuído diretamente com desenvolvimento de código nesta \textit{sprint}, mantive-me engajado no acompanhamento das tarefas do \textit{backend}. Como único \ac{ages} III atuando nesta frente, senti a responsabilidade pelos resultados, ficando frustrado por não ter conseguido atender plenamente às expectativas estabelecidas para esta etapa do projeto.
